% GENERATED FILE. Edit canonical lessons, then run npm run generate:books.
% canonical-source-hash: fnv1a64:6081f81032feb26f
% canonical-lessons: JA-C151-kimono, JA-C151-obi, JA-C151-tabi, JA-C151-sode, JA-C151-niau, JA-R151-review-146-151

\chapter{Kimono, Sashes and Sleeves}
\label{ch:ja-kimono-sashes-and-sleeves}
\hlchaptermodality{151}

\begin{chapteropening}
\textbf{By the end of this chapter:} \emph{I can name what goes with traditional dress.}

Complete the last lesson of chapter 151: I can name what goes with traditional dress.
\end{chapteropening}

\section[\texorpdfstring{\ja{きもの} (kimono)}{きもの (kimono)}]{\ja{きもの} (kimono) — a kimono}

\label{lesson:JA-C151-kimono}



\begin{quote}
Before the new one: say the Japanese for to tie, then the Japanese for a ring.
\end{quote}

\subsection*{You'll want to know: \ja{きもの}}
\textbf{\ja{きもの}} — \emph{kimono} — \textquotedblleft{}a kimono\textquotedblright{}.

\begin{quote}
\textbf{\ja{まつりで} \ja{きものを} \ja{きます。}} — \emph{matsuri de kimono o kimasu} — \textquotedblleft{}I wear a kimono at the festival.\textquotedblright{}
\end{quote}

\begin{cousinweb}
\textbf{\ja{きもの}} (\emph{kimono}) is native Japanese, built from \emph{ki}, \textquotedblleft{}wearing\textquotedblright{}, from \emph{kiru}, and \emph{mono}, \textquotedblleft{}thing\textquotedblright{}: a thing to wear. Both parts are inherited rather than borrowed (JMdict, KANJIDIC2).
\end{cousinweb}

\begin{grammarlens}[title={What you've built}]
One more word for everyday life at A2.
\end{grammarlens}

\subsection*{Your turn}
\begin{itemize}
\raggedright
  \item \emph{Say it:} \emph{kimono}
  \item \emph{Say it:} \emph{kimono}, once more
  \item \emph{Say it:} \emph{matsuri de kimono o kimasu} — I wear a kimono at the festival.
  \item \emph{Recall:} say the Japanese for to tie, then the Japanese for a ring, then say \emph{kimono} again
\end{itemize}

\begin{tcolorbox}[breakable,colback=teal!4,colframe=teal!35!black,title={Before you move on}]
What does \ja{きもの} mean? (\textquotedblleft{}A kimono\textquotedblright{}.) Say it once more.
\end{tcolorbox}
\section[\texorpdfstring{\ja{おび} (obi)}{おび (obi)}]{\ja{おび} (obi) — a kimono sash}

\label{lesson:JA-C151-obi}



\begin{quote}
Before the new one: say the Japanese for a ring, then the Japanese for a kimono.
\end{quote}

\subsection*{You'll want to know: \ja{おび}}
\textbf{\ja{おび}} — \emph{obi} — \textquotedblleft{}a kimono sash\textquotedblright{}.

\begin{quote}
\textbf{\ja{あねは} \ja{あかい} \ja{おびを} \ja{むすびます。}} — \emph{ane wa akai obi o musubimasu} — \textquotedblleft{}My elder sister ties a red sash.\textquotedblright{}
\end{quote}

\begin{cousinweb}
\textbf{\ja{おび}} (\emph{obi}) is a native Japanese word (wago), inherited rather than borrowed from Chinese; its kanji was attached later as a kun'yomi, and its older origin is not settled (JMdict, KANJIDIC2).
\end{cousinweb}

\begin{grammarlens}[title={What you've built}]
One more word for everyday life at A2.
\end{grammarlens}

\subsection*{Your turn}
\begin{itemize}
\raggedright
  \item \emph{Say it:} \emph{obi}
  \item \emph{Say it:} \emph{obi}, once more
  \item \emph{Say it:} \emph{ane wa akai obi o musubimasu} — My elder sister ties a red sash.
  \item \emph{Recall:} say the Japanese for a ring, then the Japanese for a kimono, then say \emph{obi} again
\end{itemize}

\begin{tcolorbox}[breakable,colback=teal!4,colframe=teal!35!black,title={Before you move on}]
What does \ja{おび} mean? (\textquotedblleft{}A kimono sash\textquotedblright{}.) Say it once more.
\end{tcolorbox}
\section[\texorpdfstring{\ja{たび} (tabi)}{たび (tabi)}]{\ja{たび} (tabi) — split-toe socks}

\label{lesson:JA-C151-tabi}



\begin{quote}
Before the new one: say the Japanese for a kimono, then the Japanese for a kimono sash.
\end{quote}

\subsection*{You'll want to know: \ja{たび}}
\textbf{\ja{たび}} — \emph{tabi} — \textquotedblleft{}split-toe socks\textquotedblright{}.

\begin{quote}
\textbf{\ja{まつりで} \ja{しろい} \ja{たびを} \ja{はきます。}} — \emph{matsuri de shiroi tabi o hakimasu} — \textquotedblleft{}I wear white split-toe socks at the festival.\textquotedblright{}
\end{quote}

\begin{cousinweb}
\textbf{\ja{たび}} (\emph{tabi}) is a native Japanese word, inherited rather than borrowed from Chinese, and its older origin is not settled. Its two characters, \textquotedblleft{}foot\textquotedblright{} and \textquotedblleft{}bag\textquotedblright{}, are read together as one word rather than one by one (JMdict, KANJIDIC2).
\end{cousinweb}

\begin{grammarlens}[title={What you've built}]
One more word for everyday life at A2.
\end{grammarlens}

\subsection*{Your turn}
\begin{itemize}
\raggedright
  \item \emph{Say it:} \emph{tabi}
  \item \emph{Say it:} \emph{tabi}, once more
  \item \emph{Say it:} \emph{matsuri de shiroi tabi o hakimasu} — I wear white split-toe socks at the festival.
  \item \emph{Recall:} say the Japanese for a kimono, then the Japanese for a kimono sash, then say \emph{tabi} again
\end{itemize}

\begin{tcolorbox}[breakable,colback=teal!4,colframe=teal!35!black,title={Before you move on}]
What does \ja{たび} mean? (\textquotedblleft{}Split-toe socks\textquotedblright{}.) Say it once more.
\end{tcolorbox}
\section[\texorpdfstring{\ja{そで} (sode)}{そで (sode)}]{\ja{そで} (sode) — a sleeve}

\label{lesson:JA-C151-sode}



\begin{quote}
Before the new one: say the Japanese for a kimono sash, then the Japanese for split-toe socks.
\end{quote}

\subsection*{You'll want to know: \ja{そで}}
\textbf{\ja{そで}} — \emph{sode} — \textquotedblleft{}a sleeve\textquotedblright{}.

\begin{quote}
\textbf{\ja{この} \ja{ふくは} \ja{そでが} \ja{ながいです。}} — \emph{kono fuku wa sode ga nagai desu} — \textquotedblleft{}The sleeves of these clothes are long.\textquotedblright{}
\end{quote}

\begin{cousinweb}
\textbf{\ja{そで}} (\emph{sode}) is a native Japanese word (wago), inherited rather than borrowed from Chinese; its kanji was attached later as a kun'yomi, and its older origin is not settled (JMdict, KANJIDIC2).
\end{cousinweb}

\begin{grammarlens}[title={What you've built}]
One more word for everyday life at A2.
\end{grammarlens}

\subsection*{Your turn}
\begin{itemize}
\raggedright
  \item \emph{Say it:} \emph{sode}
  \item \emph{Say it:} \emph{sode}, once more
  \item \emph{Say it:} \emph{kono fuku wa sode ga nagai desu} — The sleeves of these clothes are long.
  \item \emph{Recall:} say the Japanese for a kimono sash, then the Japanese for split-toe socks, then say \emph{sode} again
\end{itemize}

\begin{tcolorbox}[breakable,colback=teal!4,colframe=teal!35!black,title={Before you move on}]
What does \ja{そで} mean? (\textquotedblleft{}A sleeve\textquotedblright{}.) Say it once more.
\end{tcolorbox}
\section[\texorpdfstring{\ja{にあう} (niau)}{にあう (niau)}]{\ja{にあう} (niau) — to suit; to look good on}

\label{lesson:JA-C151-niau}



\begin{quote}
Before the new one: say the Japanese for split-toe socks, then the Japanese for a sleeve.
\end{quote}

\subsection*{You'll want to know: \ja{にあう}}
\textbf{\ja{にあう}} — \emph{niau} — \textquotedblleft{}to suit; to look good on\textquotedblright{}.

\begin{quote}
\textbf{\ja{この} \ja{きものは} \ja{ははに} \ja{にあいます。}} — \emph{kono kimono wa haha ni niaimasu} — \textquotedblleft{}This kimono suits my mother.\textquotedblright{}
\end{quote}

Polite \textbf{\ja{にあいます}}, \emph{niaimasu}; a godan verb.

\begin{cousinweb}
\textbf{\ja{にあう}} (\emph{niau}) is native Japanese, built from \emph{ni}, the stem of \emph{niru}, \textquotedblleft{}to resemble\textquotedblright{}, and \emph{au}, \textquotedblleft{}to match\textquotedblright{}, both inherited rather than borrowed; its kanji were attached later as kun'yomi (JMdict, KANJIDIC2).
\end{cousinweb}

\begin{grammarlens}[title={What you've built}]
One more word for everyday life at A2.
\end{grammarlens}

\subsection*{Your turn}
\begin{itemize}
\raggedright
  \item \emph{Say it:} \emph{niau}
  \item \emph{Say it:} \emph{niau}, once more
  \item \emph{Say it:} \emph{kono kimono wa haha ni niaimasu} — This kimono suits my mother.
  \item \emph{Recall:} say the Japanese for split-toe socks, then the Japanese for a sleeve, then say \emph{niau} again
\end{itemize}

\begin{tcolorbox}[breakable,colback=teal!4,colframe=teal!35!black,title={Before you move on}]
What does \ja{にあう} mean? (\textquotedblleft{}To suit; to look good on\textquotedblright{}.) Say it once more.
\end{tcolorbox}
\section[(dialogue)]{Words for the day and the home: a review}

\label{lesson:JA-R151-review-146-151}



\begin{quote}
Say the Japanese for to fall asleep and soundly.
\end{quote}

\subsection*{Your turn}
Say these aloud:

\begin{itemize}
\raggedright
  \item to fall asleep, soundly, to wake up, to wake someone, sleep
  \item to wash, to polish, to pour over oneself, hot water, soap
  \item the hair on one's head, to trim, a beard, to grow a beard, make-up
  \item to wear, to take off, to wear on the head, Western clothes, underwear
  \item a jacket, gloves, to remove, to tie, a ring
  \item a kimono, a kimono sash, split-toe socks, a sleeve, to suit
\end{itemize}

\begin{tcolorbox}[breakable,colback=teal!4,colframe=teal!35!black,title={Before you move on}]
To fall asleep? (\textbf{\ja{ねむる}}, \emph{nemuru}.) Soundly? (\textbf{\ja{ぐっすり}}, \emph{gussuri}.) To wake up? (\textbf{\ja{さめる}}, \emph{sameru}.) To wake someone? (\textbf{\ja{おこす}}, \emph{okosu}.) Sleep? (\textbf{\ja{すいみん}}, \emph{suimin}.) To wash? (\textbf{\ja{あらう}}, \emph{arau}.) To polish? (\textbf{\ja{みがく}}, \emph{migaku}.) To pour over oneself? (\textbf{\ja{あびる}}, \emph{abiru}.) Hot water? (\textbf{\ja{ゆ}}, \emph{yu}.) Soap? (\textbf{\ja{せっけん}}, \emph{sekken}.) The hair on one's head? (\textbf{\ja{かみのけ}}, \emph{kaminoke}.) To trim? (\textbf{\ja{かる}}, \emph{karu}.) A beard? (\textbf{\ja{ひげ}}, \emph{hige}.) To grow a beard? (\textbf{\ja{はやす}}, \emph{hayasu}.) Make-up? (\textbf{\ja{けしょう}}, \emph{keshō}.) To wear? (\textbf{\ja{きる}}, \emph{kiru}.) To take off? (\textbf{\ja{ぬぐ}}, \emph{nugu}.) To wear on the head? (\textbf{\ja{かぶる}}, \emph{kaburu}.) Western clothes? (\textbf{\ja{ようふく}}, \emph{yōfuku}.) Underwear? (\textbf{\ja{したぎ}}, \emph{shitagi}.) A jacket? (\textbf{\ja{うわぎ}}, \emph{uwagi}.) Gloves? (\textbf{\ja{てぶくろ}}, \emph{tebukuro}.) To remove? (\textbf{\ja{はずす}}, \emph{hazusu}.) To tie? (\textbf{\ja{むすぶ}}, \emph{musubu}.) A ring? (\textbf{\ja{ゆびわ}}, \emph{yubiwa}.) A kimono? (\textbf{\ja{きもの}}, \emph{kimono}.) A kimono sash? (\textbf{\ja{おび}}, \emph{obi}.) Split-toe socks? (\textbf{\ja{たび}}, \emph{tabi}.) A sleeve? (\textbf{\ja{そで}}, \emph{sode}.) To suit? (\textbf{\ja{にあう}}, \emph{niau}.)
\end{tcolorbox}
